\subsection{基本邻域系；拓扑基}

定义 5. 在拓扑空间 X 中，点 x（相应地，X 的子集 A）的基本邻域系是指 x（相应地，A）的邻域所成的任一集合 S，使得对 x（相应地，A）的每个邻域 V，存在邻域 W ∈ S 满足 W ⊂ V。

如果 \(\mathfrak{S}\) 是子集 \(A\)（\(X\) 的子集）的基本邻域系，则 \(\mathfrak{S}\) 中集合的每个有限交都包含 \(\mathfrak{S}\) 的某个集合。

例。1) 在离散空间（第 1 目）中，单独的集合 \(\{x\}\) 就构成点 \(x\) 的基本邻域系。

\begin{enumerate}
\def\labelenumi{\arabic{enumi})}
\setcounter{enumi}{1}
\tightlist
\item
  在有理直线 Q 上，包含点 x 的所有开区间所成之集是该点的一个基本邻域系。开区间 \(]x - 1/n, x + 1/n[\) 所成之集，以及闭区间 \([x - 1/n, x + 1/n]\) 所成之集亦然，其中 n 取遍所有整数 \textgreater0，或取遍由整数 \textgreater0 组成的任一严格递增无穷序列。* 对实直线有类似的结果。*
\end{enumerate}

定义 6. 拓扑空间 X 的拓扑的拓扑基是指 X 的开子集所成的任一集合 B，使得 X 的每个开子集都是属于 B 的集合之并。

命题 3. 如果 X 是拓扑空间，则 X 的开子集所成的集合 B 是 X 的拓扑的拓扑基的充分必要条件是：对每个 \(x \in X\)，由 \(V \in B\)（其中 \(x \in V\)）所成之集是 x 的基本邻域系。

显然该条件是必要的。反之，若该条件成立，则对任一开集 U 与任一 \(x \in U\)，存在开集 \(V_{x} \in B\) 使得 \(x \in V_{x} \subset U\)。于是诸集合 \(V_{x}\)（对 \(x \in U\)）之并等于 U。证毕。

例。1) 离散拓扑以 X 的由单独一点组成的子集所成之集为一个拓扑基。

\begin{enumerate}
\def\labelenumi{\arabic{enumi})}
\setcounter{enumi}{1}
\tightlist
\item
  有界开区间所成之集依定义是有理直线（第 2 目）的拓扑的一个拓扑基。* 同样，有界开区间所成之集是实直线的拓扑的一个拓扑基。*
\end{enumerate}

